Question answering over knowledge bases (KBQA) aims to provide a user-friendly way to access large-scale knowledge bases (KBs) by natural language questions. Existing KBQA methods~\cite{ZHANG20231} can be roughly categorized into retrieval-based and semantic-parsing (SP) based methods. The former ~\cite{feng-etal-2021-pretraining-numerical,he2021improving,zhang2022subgraph} directly scores the relevance between the question and answer candidates, thus
it is difficult to resolve the complex questions. On the contrary, some KBQA approaches, such as \cite{das2021case,kapanipathi2021leveraging,qiu2020stepwise,sun2020sparqa}, are based on semantic parsing (denoted as SP-based), which can address complex questions and achieve promising results on i.i.d. datasets. SP-based methods first translate the questions into logical expressions such as SPARQL and then execute them against KB to yield answers. As illustrated in Figure~\ref{fig:dataset}, a logical expression consists of multiple components such as classes and relations. Most existing SP-based approaches fail with logical expressions that contain unseen compositions of components (called compositional generalization) or unseen components (called zero-shot generalization).

\begin{figure}[t]
\centering
\includegraphics[scale=0.33]{dataset.png}
\caption{Illustration of generalization tasks in KBQA. Each question is paired with a logical expression that consists of different components. Components involved in the training data are colored in non-green color, while unseen components are colored in green.}     
\label{fig:dataset}

\end{figure}

\begin{figure}[t]
\centering
\includegraphics[scale=0.38]{pre_experiment.png}
\caption{Results of the pilot study. The coarse-grained method directly matches the question with the logical expression (i.e., the composition of components), while the fine-grained method matches the question with each component candidate and then composes them to derive the logical expression. The exact match accuracy of logical expressions on compositional generalization test data and zero-shot generalization test data is shown on the right of the figure.}     
\label{fig:pre_experiment}

\end{figure}

To address the above problem, GrailQA-Rank~\cite{gu2021beyond} proposes a BERT-based rank model to match the given question with each logical expression candidate, which leverages the generalization abilities of the pre-trained language models. 
On top of that, RNG-KBQA~\cite{ye2022rng} further uses a pre-trained generation model, which takes top-5 ranked logical expressions as the additional input beyond the question to generate the target logical expression.  
Behind these mainstream models, a logical expression is viewed as an inseparable unit during modeling. Actually, logical expressions are coarse-grained because they can be decomposed into relatively fine-grained components including relations, classes, entities, and logical skeletons (See examples in Figure~\ref{fig:KB}). 
Such coarse-grained modeling entangles representations of fine-grained components, thereby overfitting the seen compositions during the training process, which weakens the model's compositional generalization ability. %Meanwhile, in the zero-shot inference, such entanglement contaminates the semantic representation of each component, which weakens the zero-shot generalization ability that comes from the pre-trained language models. 
Meanwhile, even though pre-trained language models can deal with zero-shot components to some extent, compositional overfit reduces their ability to identify individual unseen components with zero-shot generalization. 

%In light of this, we explore another way to improve the generalization ability of KBQA models by decomposing the logical expression and modeling the fine-grained knowledge in it.
To demonstrate the above idea, we perform a pilot study (Cf. the detailed settings in Section~\ref{sec:pre-experiment}) with two preliminary experiments: one calculates the similarity score between a question and each coarse-grained logical expression to obtain the most relevant one, and the other searches the most relevant fine-grained components to form the final logical expression of a question. We observe that \textbf{the fine-grained modeling derives more accurate logical expressions on both the compositional task and zero-shot task (Cf. Figure~\ref{fig:pre_experiment})}. It could be explained that fine-grained modeling focuses exclusively on each component, avoiding overfitting of seen compositions in the training data. 
Although some studies attempt to leverage fine-grained components, they only consider partial fine-grained components such as relations, classes, and entities ~\cite{chen-etal-2021-retrack}, or suffer from inexecutability due to disconnected fine-grained components in real KBs~\cite{shu2022tiara}.

Thus, to both ensure the generalization ability and executability of logical expressions, we propose a \underline{F}ine-to-\underline{C}oarse composition framework for KBQA (\model), which contains three sub-modules. The overview of our model is shown in Figure~\ref{fig:over-view}.
The first module is fine-grained component detection, which detects all kinds of fine-grained component candidates from Freebase by their semantic similarities with the question. Such component detection guarantees the generalization ability in both compositional and zero-shot tasks. The second module is the middle-grained component constraint, which efficiently prunes and composes the fine-grained component candidates by ensuring the components' connectivity in the KB. %so as to improve the executability of the subsequent generated logical expression. 
The final module is the coarse-grained component composition, which employs a seq-to-seq generation model to generate the executable coarse-grained logical expression. In addition to encode the fine-grained components, the middle-grained components are also encoded to enhance the model's reasoning capacity, so as to improve the executability of the generated logical expression. In contrast to previous work~\cite{cao-etal-2022-program,chen-etal-2021-retrack,shu2022tiara} that only uses the knowledge constraints to guide the decoding process, we emphasize injecting them into the encoding process, because the encoder which learns bidirectional context could better suit natural language understanding~\cite{du2022glm}.

We conduct extensive experiments on widely used GrailQA, WebQSP, and CWQ datasets. GrailQA~\cite{gu2021beyond} is a KBQA benchmark focusing on generalization problems. \smodel derives new state-of-the-art performance on GrailQA-Dev (+7.6\% F1 gain and +7.0\% EM gain respectively). Meanwhile, \smodel also obtains good performance on WebQSP and CWQ. Moreover, \smodel runs 4 times faster than the state-of-the-art baseline RNG-KBQA. The ablation studies demonstrate the effect of our middle-grained encoding strategy.

\vpara{Contributions.}
(1) We conduct a pilot study to reveal an intriguing phenomenon --- a fine-grained understanding of the logical expression helps enhance the generalization ability of SP-based KBQA methods, which is rarely discussed before.
(2) We propose a fine-to-coarse composition framework \smodel to address the generalization problem, which takes advantage of the idea of fine-grained modeling.
(3) We devise a middle-grained component constraint that is injected into both the encoder and the decoder to guide the seq-to-seq model in producing executable logical expressions.
(4) \smodel not only maintains efficiency but also achieves significant improvement on GrailQA.